% GENERATED FILE. Edit canonical lessons, then run npm run generate:books.
% canonical-source-hash: fnv1a64:cf053b484d174152
% canonical-lessons: ZH-C293-zazhi, ZH-C293-bao, ZH-W293-bao-report, ZH-C293-jianbao, ZH-C293-toutiao, ZH-C293-guanzhu

\chapter{Headlines}
\label{ch:zh-headlines}
\hlchaptermodality{293}

\begin{chapteropening}
\textbf{By the end of this chapter:} \emph{I can talk about newspapers and magazines, the lead story and the news I follow.}

Complete the last lesson of chapter 293: I can talk about newspapers and magazines, the lead story and the news I follow.
\end{chapteropening}

\section[zázhì]{zázhì — magazine}

\label{lesson:ZH-C293-zazhi}



\begin{quote}
Before the new one: say the Mandarin for the people listening to the radio, then the Mandarin for to publish an article.
\end{quote}

\subsection*{You'll want to know: zázhì}
\emph{zázhì} — \textquotedblleft{}magazine\textquotedblright{}.

\begin{quote}
\emph{wǒ mǎi le liǎng běn zázhì} — I bought two magazines.
\end{quote}

Its characters come later; for now it is a word for the ear.

\begin{cousinweb}
\emph{zázhì} is built from \emph{zá} 'mixed' and \emph{zhì} 'a record', both inherited from Middle Chinese, \emph{dzop} and \emph{tsyiH} (Guangyun, in Baxter's transcription), which glosses the first 'to gather' and 'five colours together' and the second 'to note down': a record of many things.
\end{cousinweb}

\begin{grammarlens}[title={What you've built}]
One more word for telling what happened.
\end{grammarlens}

\subsection*{Your turn}
\begin{itemize}
\raggedright
  \item \emph{Say it:} \emph{zázhì}
  \item \emph{Say it:} \emph{zázhì}, once more
  \item \emph{Say it:} \emph{wǒ mǎi le liǎng běn zázhì} — I bought two magazines.
  \item \emph{Recall:} say the Mandarin for the people listening to the radio, then the Mandarin for to publish an article, then say \emph{zázhì} again
\end{itemize}

\begin{tcolorbox}[breakable,colback=teal!4,colframe=teal!35!black,title={Before you move on}]
What does zázhì mean? (\textquotedblleft{}Magazine\textquotedblright{}.) Say it once more.
\end{tcolorbox}
\section[\texorpdfstring{\zh{报} (bào)}{报 (bào)}]{\zh{报} (bào) — newspaper; the daily paper}

\label{lesson:ZH-C293-bao}



\begin{quote}
Before the new one: say the Mandarin for to publish an article, then the Mandarin for a magazine.
\end{quote}

\subsection*{You'll want to know: \zh{报}}
\textbf{\zh{报}} — \emph{bào} — \textquotedblleft{}newspaper; the daily paper\textquotedblright{}.

\begin{quote}
\textbf{\zh{我每天早上看报}} \emph{wǒ měitiān zǎoshang kàn bào} — I read the paper every morning.
\end{quote}

The new character, in pieces: \zh{报} — seven strokes in two pieces: on the left the hand piece of \zh{打}, three strokes; on the right a bar that turns down and hooks back, a long vertical down its left side, then a bar that turns into a sweep to the lower left and a long stroke falling to the right. The next lesson puts it in your hand.

\begin{cousinweb}
\textbf{\zh{报}} (\emph{bào}) is inherited from Middle Chinese \emph{pawH} (Guangyun, in Baxter's transcription), which glosses it 'to report'; the sense 'a newspaper', the sheet that reports, is later and is not given there.
\end{cousinweb}

\begin{grammarlens}[title={What you've built}]
One more word for telling what happened.
\end{grammarlens}

\subsection*{Your turn}
\begin{itemize}
\raggedright
  \item \emph{Say it:} \emph{bào}
  \item \emph{Say it:} \emph{bào}, once more
  \item \emph{Say it:} \emph{wǒ měitiān zǎoshang kàn bào} — I read the paper every morning.
  \item \emph{Recall:} say the Mandarin for to publish an article, then the Mandarin for a magazine, then say \emph{bào} again
\end{itemize}

\begin{tcolorbox}[breakable,colback=teal!4,colframe=teal!35!black,title={Before you move on}]
What does \zh{报} mean? (\textquotedblleft{}Newspaper; the daily paper\textquotedblright{}.) Say it once more.
\end{tcolorbox}
\section[\texorpdfstring{\zh{报} (bào)}{报 (bào)}]{\zh{报} — the hand piece, then a hook, a post and a fork}

\label{lesson:ZH-W293-bao-report}



\begin{quote}
The word first, then the character inside it.

\begin{itemize}
\raggedright
  \item \emph{Recall:} say \emph{bào} — \textbf{R1}, one lesson back, before it decays
\end{itemize}

\begin{quote}
\zh{打} \zh{找}
\end{quote}

\emph{Read it:} each one aloud, with nothing said first — \emph{dǎ}, \emph{zhǎo}
\end{quote}

\subsection*{Script — the hand piece, then a hook, a post and a fork}
\hlblockfigure{\detokenize{figures/ZH-W293-bao-report-filmstrip.pdf}}{How \zh{报} is written, stroke by stroke}

\begin{quote}
\zh{报}
\end{quote}

\textbf{Two pieces, left before right.} Count the pieces, not the strokes.

\textbf{The left piece, three strokes.} The hand piece of \zh{打}: a bar, a vertical that hooks to the left at the foot, and a stroke rising to the right across it.

\textbf{The right piece, four strokes.} A bar that turns down and hooks back to the left. A long vertical down the left side of the piece. Then the fork: a bar that turns into a sweep down to the lower left, and a long stroke falling to the right.

3 + 4 = seven.

\subsection*{Your turn}
\begin{itemize}
\raggedright
  \item \emph{Write it:} \zh{报}, all seven, with the model in front of you
  \item \emph{Cover:} write it again from memory, and count the pieces, not the strokes
\end{itemize}

\begin{tcolorbox}[breakable,colback=teal!4,colframe=teal!35!black,title={Before you move on}]
How many pieces in \zh{报}? (\textbf{Two}.) What is the left piece? (\textbf{The hand piece of \zh{打}, three strokes}.) What closes the right piece? (\textbf{A fork: a sweep to the left, then a long stroke to the right}.)

Source: \href{https://raw.githubusercontent.com/chanind/hanzi-writer-data/68d10a4b21150cae5e1ebbd223eed289cf32d90c/data/\%E6\%8A\%A5.json}{Hanzi Writer Data \zh{报}, pinned stroke paths and medians}.
\end{tcolorbox}
\section[\texorpdfstring{\zh{见报} (jiànbào)}{见报 (jiànbào)}]{\zh{见报} (jiànbào) — to appear in the papers; to make the papers}

\label{lesson:ZH-C293-jianbao}



\begin{quote}
Before the new one: say the Mandarin for a magazine, then the Mandarin for the daily paper.
\end{quote}

\subsection*{You'll want to know: \zh{见报}}
\textbf{\zh{见报}} — \emph{jiànbào} — \textquotedblleft{}to appear in the papers; to make the papers\textquotedblright{}.

\begin{quote}
\textbf{\zh{我们学校见报了}} \emph{wǒmen xuéxiào jiànbào le} — Our school was in the paper.
\end{quote}

\begin{cousinweb}
\textbf{\zh{见报}} (\emph{jiànbào}) is built from \zh{见} 'to see' and \zh{报} 'a newspaper', both inherited from Middle Chinese, \emph{kenH} and \emph{pawH} (Guangyun, in Baxter's transcription), which glosses the first 'to view': to be seen in the paper.
\end{cousinweb}

\begin{grammarlens}[title={What you've built}]
One more word for telling what happened.
\end{grammarlens}

\subsection*{Your turn}
\begin{itemize}
\raggedright
  \item \emph{Say it:} \emph{jiànbào}
  \item \emph{Say it:} \emph{jiànbào}, once more
  \item \emph{Say it:} \emph{wǒmen xuéxiào jiànbào le} — Our school was in the paper.
  \item \emph{Recall:} say the Mandarin for a magazine, then the Mandarin for the daily paper, then say \emph{jiànbào} again
\end{itemize}

\begin{tcolorbox}[breakable,colback=teal!4,colframe=teal!35!black,title={Before you move on}]
What does \zh{见报} mean? (\textquotedblleft{}To appear in the papers; to make the papers\textquotedblright{}.) Say it once more.
\end{tcolorbox}
\section[tóutiáo]{tóutiáo — lead story; headline; the front-page story}

\label{lesson:ZH-C293-toutiao}



\begin{quote}
Before the new one: say the Mandarin for the daily paper, then the Mandarin for to make the papers.
\end{quote}

\subsection*{You'll want to know: tóutiáo}
\emph{tóutiáo} — \textquotedblleft{}lead story; headline; the front-page story\textquotedblright{}.

\begin{quote}
\emph{jīntiān de tóutiáo shì shénme} — What's today's lead story?
\end{quote}

Its characters come later; for now it is a word for the ear.

\begin{cousinweb}
\emph{tóutiáo} is built from \zh{头} 'the head' and \emph{tiáo} 'a strip', both inherited from Middle Chinese, \emph{duw} and \emph{dew} (Guangyun, in Baxter's transcription), which glosses the first 'the head' and the second 'a small branch': the first strip, the story at the top.
\end{cousinweb}

\begin{grammarlens}[title={What you've built}]
One more word for telling what happened.
\end{grammarlens}

\subsection*{Your turn}
\begin{itemize}
\raggedright
  \item \emph{Say it:} \emph{tóutiáo}
  \item \emph{Say it:} \emph{tóutiáo}, once more
  \item \emph{Say it:} \emph{jīntiān de tóutiáo shì shénme} — What's today's lead story?
  \item \emph{Recall:} say the Mandarin for the daily paper, then the Mandarin for to make the papers, then say \emph{tóutiáo} again
\end{itemize}

\begin{tcolorbox}[breakable,colback=teal!4,colframe=teal!35!black,title={Before you move on}]
What does tóutiáo mean? (\textquotedblleft{}Lead story; headline; the front-page story\textquotedblright{}.) Say it once more.
\end{tcolorbox}
\section[guānzhù]{guānzhù — to follow a story closely; to keep an eye on}

\label{lesson:ZH-C293-guanzhu}



\begin{quote}
Before the new one: say the Mandarin for to make the papers, then the Mandarin for the front-page story.
\end{quote}

\subsection*{You'll want to know: guānzhù}
\emph{guānzhù} — \textquotedblleft{}to follow a story closely; to keep an eye on\textquotedblright{}.

\begin{quote}
\emph{rúguǒ nǐ guānzhù xīnwén, nǐ jiù huì zhīdào} — If you followed the news, you'd know.
\end{quote}

Its characters come later; for now it is a word for the ear.

\begin{cousinweb}
\emph{guānzhù} is built from \zh{关} 'to shut' and \emph{zhù} 'to pour', both inherited from Middle Chinese, \emph{kwaen} and \emph{tsyuH} (Guangyun, in Baxter's transcription), which glosses the first 'the bar that holds a gate shut' and the second 'to pour into'; how the pair came to mean 'to follow closely' is not settled.
\end{cousinweb}

\begin{grammarlens}[title={What you've built}]
One more word for telling what happened.
\end{grammarlens}

\subsection*{Your turn}
\begin{itemize}
\raggedright
  \item \emph{Say it:} \emph{guānzhù}
  \item \emph{Say it:} \emph{guānzhù}, once more
  \item \emph{Say it:} \emph{rúguǒ nǐ guānzhù xīnwén, nǐ jiù huì zhīdào} — If you followed the news, you'd know.
  \item \emph{Recall:} say the Mandarin for to make the papers, then the Mandarin for the front-page story, then say \emph{guānzhù} again
\end{itemize}

\begin{tcolorbox}[breakable,colback=teal!4,colframe=teal!35!black,title={Before you move on}]
What does guānzhù mean? (\textquotedblleft{}To follow a story closely; to keep an eye on\textquotedblright{}.) Say it once more.
\end{tcolorbox}
